% ============================================================
% UNIDAD: SISTEMA IMPOSITIVO ARGENTINO, MONOTRIBUTO, FACTURACIÓN
% Y FINANCIAMIENTO BANCARIO
% Datos normativos y montos verificados en ARCA y BCRA el 25/11/2026.
% Los montos del monotributo se actualizan en febrero y agosto de cada año.
% ============================================================
\clearpage
\section*{UNIDAD: SISTEMA IMPOSITIVO ARGENTINO}
\subsection*{Monotributo, facturación y financiamiento para emprendedores}

\begin{recuadro}
\textbf{Importante:} los montos de esta unidad (categorías del monotributo, cuotas y topes)
son los vigentes desde el \textbf{1 de agosto de 2026} hasta enero de 2027, según la tabla
publicada por ARCA. Se actualizan dos veces por año (febrero y agosto) según la inflación,
por lo que siempre hay que verificarlos en \textbf{arca.gob.ar/monotributo}.
\end{recuadro}

% ------------------------------------------------------------
\section*{1. LOS IMPUESTOS Y QUIÉN LOS COBRA}

Un \textbf{impuesto} es una suma de dinero que el Estado exige a las personas y empresas,
sin una contraprestación directa, para financiar el gasto público (salud, educación,
seguridad, obras). En la Argentina hay tres niveles de gobierno y cada uno cobra sus propios
tributos:

\begin{center}
\small
\begin{tabular}{|p{0.17\linewidth}|p{0.22\linewidth}|p{0.51\linewidth}|}
\hline
\textbf{Nivel} & \textbf{Organismo} & \textbf{Principales tributos} \\ \hline
Nacional & ARCA (Agencia de Recaudación y Control Aduanero, ex AFIP) &
IVA, Impuesto a las Ganancias, Bienes Personales, impuestos internos, aportes y
contribuciones de la seguridad social, derechos de aduana, Monotributo. \\ \hline
Provincial & ARBA en la provincia de Buenos Aires (AGIP en CABA) &
Ingresos Brutos (IIBB), Impuesto Inmobiliario, Automotor, Sellos. \\ \hline
Municipal & Municipio (por ejemplo, Quilmes) &
Tasas: Seguridad e Higiene (comercios e industrias), alumbrado, barrido y limpieza (ABL),
publicidad y propaganda, habilitaciones. \\ \hline
\end{tabular}
\end{center}

\paragraph*{Impuestos directos e indirectos.}
Los \textbf{directos} gravan la riqueza o el ingreso de una persona (Ganancias, Bienes
Personales). Los \textbf{indirectos} gravan el consumo y se trasladan al precio que paga el
consumidor final (IVA, impuestos internos, Ingresos Brutos).

\paragraph*{Las dos formas de estar inscripto ante ARCA.}
Una persona que trabaja por su cuenta o tiene un emprendimiento puede estar en:
\begin{itemize}
  \item \textbf{Régimen Simplificado (Monotributo):} para pequeños contribuyentes. Se paga
        una cuota fija mensual que reemplaza al IVA, a Ganancias y a los aportes jubilatorios
        y de obra social.
  \item \textbf{Régimen General (Responsable Inscripto):} para quienes superan los límites
        del monotributo o no cumplen sus requisitos. Liquida IVA todos los meses, Ganancias
        una vez por año y aporta como trabajador autónomo.
\end{itemize}

% ------------------------------------------------------------
\section*{2. CUIT Y CLAVE FISCAL}

\paragraph*{CUIT.}
La \textbf{Clave Única de Identificación Tributaria} es el número con el que ARCA identifica
a cada contribuyente. Tiene 11 dígitos con el formato \textbf{XX-XXXXXXXX-X}:
\begin{itemize}
  \item \textbf{Prefijo (2 dígitos):} indica el tipo de persona. Para personas humanas se usan
        20, 23, 24 y 27; para personas jurídicas (sociedades, cooperativas, asociaciones)
        30, 33 y 34.
  \item \textbf{Número central (8 dígitos):} en las personas humanas es el número de DNI.
  \item \textbf{Dígito verificador (1 dígito):} se calcula a partir de los anteriores y sirve
        para detectar errores de tipeo.
\end{itemize}
Ejemplo: 20-12345678-9. El CUIL que figura en el recibo de sueldo tiene el mismo formato,
pero identifica a la persona ante la seguridad social (ANSES); al inscribirse en un impuesto,
la persona pasa a usar CUIT.

\paragraph*{Clave fiscal.}
Es la contraseña personal para operar los servicios web de ARCA. Tiene niveles de seguridad
del 1 al 4; cada servicio exige un nivel mínimo. Para inscribirse en el monotributo y
\textbf{emitir facturas electrónicas se necesita nivel 3 o superior}. El nivel 3 se obtiene
validando la identidad con el DNI, por ejemplo desde la aplicación móvil de ARCA o en una
dependencia del organismo.

\paragraph*{Domicilio fiscal electrónico.}
Es el lugar donde ARCA envía notificaciones oficiales (intimaciones, avisos). Es obligatorio
constituirlo y revisarlo periódicamente.

% ------------------------------------------------------------
\section*{3. EL MONOTRIBUTO}

\begin{recuadro}
El \textbf{Monotributo} (Régimen Simplificado para Pequeños Contribuyentes) unifica en
\textbf{una sola cuota mensual} el componente impositivo (que reemplaza al IVA y a
Ganancias) y el componente previsional (jubilación y obra social).
\end{recuadro}

\subsubsection*{3.1. Quiénes pueden adherir}
Personas humanas que vendan cosas muebles, presten servicios o realicen obras o locaciones,
y que cumplan, entre otras, estas condiciones:
\begin{itemize}
  \item Que sus ingresos brutos de los últimos 12 meses no superen el tope de la categoría
        más alta (K).
  \item Que no superen los máximos de superficie afectada, energía eléctrica consumida y
        alquileres de la categoría más alta.
  \item Si venden cosas muebles, que el precio unitario de venta no supere el máximo fijado
        (\$716.840,77 desde agosto de 2026).
  \item Que \textbf{no desarrollen más de tres actividades simultáneas ni tengan más de tres
        unidades de explotación} (por ejemplo, tres locales).
  \item Que no hayan realizado importaciones de bienes o servicios para revenderlos en los
        últimos 12 meses.
\end{itemize}
Quien no cumple estos requisitos debe inscribirse en el Régimen General.

\subsubsection*{3.2. Servicios y venta de cosas muebles}
ARCA distingue dos tipos de actividad, que pagan cuotas distintas:
\begin{itemize}
  \item \textbf{Locaciones y prestaciones de servicios:} se vende trabajo o conocimiento.
        Ejemplos: reparación de equipos de aviónica, mantenimiento eléctrico, diseño,
        programación, clases particulares.
  \item \textbf{Venta de cosas muebles:} se venden bienes materiales que se pueden trasladar.
        Ejemplos: venta de repuestos, componentes electrónicos, herramientas, indumentaria.
\end{itemize}
En las categorías A y B la cuota es la misma para ambos tipos. Desde la C en adelante, la
cuota de servicios es más alta que la de venta de bienes. Tanto quienes prestan servicios
como quienes venden bienes pueden encuadrarse en cualquier categoría, de la A a la K.
Declarar venta de productos cuando en realidad se prestan servicios es una causal de
exclusión.

\subsubsection*{3.3. Los parámetros y la categoría}
La categoría (de la A a la K) se determina por el \textbf{mayor} de estos parámetros:
ingresos brutos anuales, superficie afectada a la actividad, energía eléctrica consumida
en el año y alquileres devengados en el año. Es decir, si por ingresos corresponde la C pero
por superficie corresponde la E, la categoría es la E.

\begin{center}
\footnotesize
\textbf{Monotributo: parámetros y cuota mensual total (vigentes desde el 1/08/2026)}\\[3pt]
\setlength{\tabcolsep}{3.5pt}
\begin{tabular}{|c|r|c|c|r|r|r|}
\hline
\textbf{Cat.} & \textbf{Ingresos brutos} & \textbf{Superficie} & \textbf{Energía} &
\textbf{Alquileres} & \multicolumn{2}{c|}{\textbf{Cuota mensual total}} \\
 & \textbf{anuales hasta} & \textbf{hasta} & \textbf{kWh/año} & \textbf{anuales hasta} &
\textbf{Servicios} & \textbf{Venta bienes} \\ \hline
A & \$12.009.410,45 & 30 m² & 3.330 & \$2.792.886,15 & \$49.527,18 & \$49.527,18 \\ \hline
B & \$17.595.182,74 & 45 m² & 5.000 & \$2.792.886,15 & \$56.379,08 & \$56.379,08 \\ \hline
C & \$24.670.494,31 & 60 m² & 6.700 & \$3.816.944,41 & \$66.020,12 & \$64.530,58 \\ \hline
D & \$30.628.651,43 & 85 m² & 10.000 & \$3.816.944,41 & \$84.612,93 & \$82.564,81 \\ \hline
E & \$36.028.231,33 & 110 m² & 13.000 & \$4.841.002,66 & \$119.811,45 & \$108.267,51 \\ \hline
F & \$45.151.659,41 & 150 m² & 16.500 & \$4.841.002,66 & \$150.784,21 & \$129.930,65 \\ \hline
G & \$53.995.798,87 & 200 m² & 20.000 & \$5.771.964,69 & \$230.312,94 & \$158.815,05 \\ \hline
H & \$81.924.660,37 & 200 m² & 20.000 & \$8.378.658,45 & \$522.706,68 & \$317.895,01 \\ \hline
I & \$91.699.761,90 & 200 m² & 20.000 & \$8.378.658,45 & \$963.747,86 & \$474.992,78 \\ \hline
J & \$105.012.519,20 & 200 m² & 20.000 & \$8.378.658,45 & \$1.167.299,76 & \$580.793,69 \\ \hline
K & \$126.610.838,75 & 200 m² & 20.000 & \$8.378.658,45 & \$1.614.446,04 & \$702.103,24 \\ \hline
\end{tabular}\\[2pt]
{\scriptsize Fuente: ARCA, Montos y categorías vigentes. Precio unitario máximo para venta de
cosas muebles: \$716.840,77. El parámetro de superficie no se aplica en localidades de hasta
40.000 habitantes.}
\end{center}

\subsubsection*{3.4. Qué incluye la cuota}
La cuota mensual tiene tres componentes:
\begin{enumerate}
  \item \textbf{Impuesto integrado:} es la parte impositiva; reemplaza al IVA y a Ganancias.
        Varía según la categoría y el tipo de actividad.
  \item \textbf{Aporte al SIPA:} el \textbf{Sistema Integrado Previsional Argentino} es el
        sistema público de jubilaciones administrado por ANSES. Este aporte computa años de
        servicio para la jubilación.
  \item \textbf{Aporte a la obra social:} da cobertura de salud al titular. Por cada
        familiar adherente (cónyuge, hijos) se paga además el mismo importe.
\end{enumerate}

\begin{center}
\footnotesize
\textbf{Componentes de la cuota en algunas categorías (desde el 1/08/2026)}\\[3pt]
\begin{tabular}{|c|r|r|r|r|}
\hline
\textbf{Cat.} & \textbf{Imp. integrado servicios} & \textbf{Imp. integrado venta} &
\textbf{Aporte SIPA} & \textbf{Obra social} \\ \hline
A & \$5.585,77 & \$5.585,77 & \$18.246,86 & \$25.694,55 \\ \hline
C & \$18.246,86 & \$16.757,32 & \$22.078,71 & \$25.694,55 \\ \hline
E & \$55.857,73 & \$44.313,79 & \$26.715,24 & \$37.238,48 \\ \hline
H & \$409.623,31 & \$204.811,64 & \$57.598,04 & \$55.485,33 \\ \hline
K & \$1.368.514,20 & \$456.171,40 & \$158.049,02 & \$87.882,82 \\ \hline
\end{tabular}\\[2pt]
{\scriptsize Quien ya aporta a otro régimen previsional obligatorio (por ejemplo, un
trabajador en relación de dependencia) está exceptuado de pagar el SIPA y la obra social
del monotributo; los jubilados no pagan obra social.}
\end{center}

\subsubsection*{3.5. Obligaciones del monotributista}
\begin{itemize}
  \item \textbf{Pagar la cuota todos los meses}, con vencimiento el día 20. Acumular deuda
        de varias cuotas puede llevar a la baja automática del régimen.
  \item \textbf{Facturar todas las operaciones} (factura C).
  \item \textbf{Recategorizarse} en febrero y agosto si cambiaron sus parámetros de los
        últimos 12 meses. Se hace en el servicio ``Monotributo'' con clave fiscal.
  \item Exhibir en el local la credencial del monotributo y el comprobante del último pago.
\end{itemize}
Si se superan los límites de la categoría más alta, se vende por encima del precio unitario
máximo o se incumplen los requisitos, ARCA \textbf{excluye} al contribuyente y debe pasar al
Régimen General.

\subsubsection*{3.6. Ingresos Brutos, Monotributo Unificado y Convenio Multilateral}
El monotributo es un régimen \textbf{nacional}. Además, quien realiza una actividad
comercial debe pagar \textbf{Ingresos Brutos}, que es un impuesto \textbf{provincial}.
\begin{itemize}
  \item \textbf{Monotributo Unificado:} varias provincias firmaron convenios con ARCA para
        cobrar en una sola cuota el monotributo nacional y los Ingresos Brutos simplificados.
        En la provincia de Buenos Aires se llama \textbf{Régimen Simplificado de Ingresos
        Brutos} y lo administra ARBA; desde enero de 2026 también funciona en CABA. El
        contribuyente no presenta declaraciones juradas de IIBB: paga un importe fijo según
        su categoría junto con el monotributo.
  \item \textbf{Convenio Multilateral:} es el régimen de Ingresos Brutos para quienes
        trabajan en \textbf{más de una jurisdicción} (por ejemplo, un taller en Quilmes que
        también presta servicios en CABA y en Córdoba). Los ingresos se reparten entre las
        provincias según dónde se generan, para que no se pague dos veces el mismo impuesto.
        Quien tiene domicilio o actividad fuera de una sola provincia deja de estar en el
        régimen simplificado y debe inscribirse en el Convenio Multilateral.
\end{itemize}

% ------------------------------------------------------------
\section*{4. EL IMPUESTO AL VALOR AGREGADO (IVA)}

\begin{recuadro}
El \textbf{IVA} es un impuesto nacional e indirecto que grava el valor que se agrega en cada
etapa de la producción y comercialización de bienes y servicios. Lo termina pagando el
\textbf{consumidor final}, dentro del precio.
\end{recuadro}

\paragraph*{Cómo funciona: débito y crédito fiscal.}
Cada responsable inscripto cobra IVA en sus ventas (\textbf{débito fiscal}) y paga IVA en
sus compras (\textbf{crédito fiscal}). Todos los meses ingresa a ARCA la diferencia:
\[
\mbox{IVA a pagar} = \mbox{Débito fiscal (IVA de las ventas)} - \mbox{Crédito fiscal (IVA de las compras)}
\]
Ejemplo: un taller compra repuestos por \$100.000 + IVA 21\% (\$21.000 de crédito fiscal) y
cobra un trabajo por \$300.000 + IVA 21\% (\$63.000 de débito fiscal). Ese mes ingresa a
ARCA \$63.000 $-$ \$21.000 = \textbf{\$42.000}. Así, cada etapa paga el impuesto solo
sobre el valor que agregó.

\paragraph*{Alícuotas.}
\begin{itemize}
  \item \textbf{21\% (general):} se aplica a la mayoría de los bienes y servicios.
  \item \textbf{10,5\% (reducida):} entre otros, carnes, frutas, verduras y hortalizas
        frescas, pan y harina de trigo, la construcción de viviendas, la medicina prepaga y
        muchos bienes de capital (máquinas y equipos).
  \item \textbf{27\% (incrementada):} servicios de electricidad, gas, agua y
        telecomunicaciones que se prestan a responsables inscriptos (no a hogares).
\end{itemize}
Existen además algunas alícuotas especiales y operaciones exentas (por ejemplo, libros, o
la venta de leche fluida a consumidores finales).

\paragraph*{El monotributista y el IVA.}
El monotributista \textbf{no cobra IVA} en sus facturas (ya está incluido en el impuesto
integrado de su cuota) y \textbf{no puede descontar} el IVA de sus compras.

\paragraph*{Régimen de Transparencia Fiscal al Consumidor.}
Desde la Ley 27.743, en los tickets y facturas a consumidores finales debe figurar un
apartado con la leyenda ``Régimen de Transparencia Fiscal al Consumidor (Ley 27.743)'' que
informa el \textbf{IVA contenido} en el precio y, si corresponde, otros impuestos
nacionales indirectos. Por ejemplo, en un precio final de \$121.000 con IVA 21\%, el IVA
contenido es \$21.000.

% ------------------------------------------------------------
\section*{5. LOS COMPROBANTES: TIPOS DE FACTURA}

Toda venta o prestación debe respaldarse con un comprobante. La letra de la factura
depende de la condición frente al IVA de \textbf{quien la emite} y de \textbf{quien la
recibe}. Hoy, en casi todos los casos, la factura es \textbf{electrónica}: ARCA la autoriza
y le asigna un \textbf{CAE} (Código de Autorización Electrónico) con fecha de vencimiento.

\begin{center}
\small
\begin{tabular}{|c|p{0.25\linewidth}|p{0.25\linewidth}|p{0.30\linewidth}|}
\hline
\textbf{Letra} & \textbf{Quién la emite} & \textbf{A quién} & \textbf{Características} \\ \hline
\textbf{A} & Responsable inscripto & Otro responsable inscripto o un monotributista &
Discrimina el IVA. El comprador inscripto lo computa como crédito fiscal. \\ \hline
\textbf{B} & Responsable inscripto & Consumidor final o sujeto exento &
El precio incluye el IVA; se informa el IVA contenido (Transparencia Fiscal). \\ \hline
\textbf{C} & Monotributista o sujeto exento & Cualquier cliente &
No discrimina IVA. Es la factura del monotributista. \\ \hline
\textbf{E} & Responsable inscripto o monotributista & Cliente del exterior &
Factura de \textbf{exportación} de bienes o servicios. \\ \hline
\textbf{T} & Hoteles y alojamientos (responsables inscriptos) & Turistas extranjeros
no residentes & Permite el beneficio de reintegro del IVA por alojamiento a turistas. \\ \hline
\end{tabular}
\end{center}

\paragraph*{La desaparecida factura M y las nuevas facturas A.}
Hasta noviembre de 2025 existía la factura \textbf{M}, que ARCA asignaba a responsables
inscriptos sin solvencia demostrada o con inconsistencias fiscales; sus clientes debían
retenerles IVA y Ganancias. La \textbf{Resolución General 5762/2025} la eliminó desde el
\textbf{1 de diciembre de 2025} y la reemplazó por dos modalidades de factura A:
\begin{itemize}
  \item \textbf{A con leyenda ``OPERACIÓN SUJETA A RETENCIÓN'':} el cliente debe retener IVA
        y Ganancias al pagar, y el cobro tiene que hacerse en la cuenta bancaria (CBU)
        informada.
  \item \textbf{A con leyenda ``PAGO EN CBU INFORMADA'':} no tiene retenciones, pero solo
        puede cobrarse por transferencia o depósito en la cuenta declarada ante ARCA.
\end{itemize}

\paragraph*{Otros comprobantes.}
\begin{itemize}
  \item \textbf{Nota de crédito:} anula total o parcialmente una factura (devoluciones,
        descuentos, errores). Lleva la misma letra que la factura que corrige.
  \item \textbf{Nota de débito:} aumenta el importe de una factura ya emitida (intereses,
        gastos omitidos).
  \item \textbf{Recibo:} constancia de un cobro.
  \item \textbf{Remito:} acompaña el traslado de mercadería; no es una factura.
  \item \textbf{Ticket:} lo emiten los controladores fiscales en las ventas al público.
\end{itemize}

\paragraph*{La ventaja impositiva de la factura A.}
Para el comprador responsable inscripto, la factura A le permite \textbf{computar el IVA
como crédito fiscal} (lo descuenta del IVA que debe pagar) y deducir el gasto en
Ganancias. Por eso las empresas suelen pedir factura A a sus proveedores inscriptos.

% ------------------------------------------------------------
\section*{6. PASO A PASO: EMITIR UNA FACTURA C EN ARCA}

\subsubsection*{6.1. Antes de facturar}
\begin{enumerate}
  \item \textbf{Tener CUIT y clave fiscal nivel 3.}
  \item \textbf{Adherir al monotributo:} en arca.gob.ar, ingresar con CUIT y clave fiscal al
        servicio \textbf{``Monotributo''} y elegir la opción de adhesión. El sistema pide
        la actividad (según el nomenclador de actividades de ARCA), los ingresos estimados,
        la superficie, el consumo de energía y el alquiler; con esos datos asigna la
        categoría. También se elige la obra social. El trámite es gratuito.
  \item \textbf{Dar de alta un punto de venta:} es el número de 4 o 5 dígitos que aparece
        al principio de cada factura (en la factura 00001-00000015, el punto de venta es el
        00001). Se crea desde el servicio de administración de puntos de venta (también se
        accede desde ``Comprobantes en línea'' $\rightarrow$ ``ABM Puntos de Venta''),
        indicando como sistema \textbf{``Factura en línea – Monotributo''}. Para hacerlo, el
        domicilio tiene que estar cargado en el Sistema Registral.
\end{enumerate}

\subsubsection*{6.2. Emitir la factura en ``Comprobantes en línea''}
\begin{enumerate}
  \item Ingresar a arca.gob.ar con CUIT y clave fiscal y abrir el servicio
        \textbf{``Comprobantes en línea''}. Elegir la persona por la que se va a facturar.
  \item Seleccionar \textbf{``Generar Comprobantes''}.
  \item Elegir el \textbf{punto de venta} y el tipo de comprobante: \textbf{Factura C}.
  \item \textbf{Datos de emisión:} fecha del comprobante y \textbf{concepto} a incluir
        (productos, servicios, o productos y servicios). Si son servicios, se indica el
        período facturado (desde/hasta, dentro del mismo mes) y la fecha de vencimiento
        del pago.
  \item \textbf{Datos del receptor:} condición frente al IVA del cliente (consumidor final,
        responsable inscripto, monotributista, exento), tipo y número de documento (CUIT o
        DNI) y condición de venta (contado, cuenta corriente, tarjeta, transferencia, etc.).
        Al ingresar una CUIT, el sistema completa la razón social y el domicilio.
  \item \textbf{Datos de la operación:} descripción de cada producto o servicio, cantidad,
        unidad de medida y precio unitario. El sistema calcula el total.
  \item Revisar el resumen y presionar \textbf{``Confirmar Datos''}. ARCA autoriza la
        factura y le asigna el \textbf{CAE}.
  \item Descargar el PDF (original y duplicado) y entregarlo o enviarlo al cliente.
\end{enumerate}

\begin{recuadro}
\textbf{Alternativas:} ARCA también ofrece la aplicación móvil \textbf{``Facturador''}, para
emitir comprobantes desde el celular; en ese caso el punto de venta se crea
automáticamente en el primer uso. Los comercios con mucho movimiento usan
\textbf{controladores fiscales} o un software de gestión conectado a ARCA por
\textit{web service}.
\end{recuadro}

\paragraph*{Errores frecuentes.}
Elegir mal la condición frente al IVA del cliente; facturar un servicio sin indicar el
período; facturar en un punto de venta dado de alta para otro sistema; no controlar la
facturación acumulada (ARCA ofrece un \textbf{Monitor de Facturación} para seguirla).
Una factura con errores no se borra: se anula con una \textbf{nota de crédito}.

% ------------------------------------------------------------
\section*{7. COMERCIO EXTERIOR E INCOTERMS}

Muchos emprendimientos técnicos compran en el exterior (repuestos, instrumentos,
componentes electrónicos) o venden fuera del país. En una compraventa internacional hay
que acordar \textbf{quién paga cada tramo del transporte, quién contrata el seguro, quién
hace los trámites de aduana y en qué momento el riesgo pasa del vendedor al comprador}.
Para eso existen los Incoterms.

\begin{recuadro}
Los \textbf{Incoterms} (\textit{International Commercial Terms}) son reglas publicadas por la
\textbf{Cámara de Comercio Internacional (ICC)} que definen, con una sigla de tres letras,
las obligaciones, los costos y los riesgos del vendedor y del comprador en una operación
internacional. La versión vigente es \textbf{Incoterms 2020}, en vigor desde el 1 de enero
de 2020. No existe una versión 2025 ni 2026: la ICC recién comenzó en 2026 los trabajos
para la próxima revisión (Incoterms 2030).
\end{recuadro}

\paragraph*{Cómo se escribe un Incoterm.}
Siempre se indica la sigla, el \textbf{lugar} y la versión. Por ejemplo:
\textbf{FCA Miami, Incoterms 2020}. Un Incoterm sin un lugar preciso no sirve.

\paragraph*{Qué NO regulan.}
No fijan el precio ni la forma de pago, no transfieren la propiedad de la mercadería y no
resuelven reclamos por calidad: eso se define en el contrato de compraventa.

\subsubsection*{7.1. Las 11 reglas Incoterms 2020}
Se dividen en dos grupos según el medio de transporte. En la tabla están ordenadas de menor
a mayor obligación para el vendedor.

\begin{center}
\small
\begin{tabular}{|c|>{\raggedright\arraybackslash}p{0.22\linewidth}|p{0.60\linewidth}|}
\hline
\textbf{Sigla} & \textbf{Nombre} & \textbf{Qué significa} \\ \hline
\multicolumn{3}{|l|}{\textbf{Para cualquier medio de transporte (incluido el aéreo y el courier)}} \\ \hline
EXW & \textit{Ex Works} (en fábrica) & El vendedor deja la mercadería disponible en su
local. El comprador se encarga de todo, incluso del despacho de exportación. \\ \hline
FCA & \textit{Free Carrier} (franco transportista) & El vendedor entrega la mercadería,
despachada de exportación, al transportista que eligió el comprador, en el lugar acordado. \\ \hline
CPT & \textit{Carriage Paid To} (transporte pagado hasta) & El vendedor paga el flete hasta
el destino, pero el riesgo pasa al comprador al entregar al primer transportista. \\ \hline
CIP & \textit{Carriage and Insurance Paid To} (transporte y seguro pagados hasta) & Como CPT,
pero el vendedor también contrata un seguro de cobertura amplia. \\ \hline
DAP & \textit{Delivered at Place} (entregado en lugar) & El vendedor lleva la mercadería
hasta el lugar de destino, lista para descargar. El comprador hace la importación. \\ \hline
DPU & \textit{Delivered at Place Unloaded} (entregado en lugar descargado) & Como DAP, pero
el vendedor también descarga la mercadería en destino. Reemplazó a DAT en 2020. \\ \hline
DDP & \textit{Delivered Duty Paid} (entregado con derechos pagados) & El vendedor asume
todo: transporte, seguro, aduana de importación e impuestos en destino. \\ \hline
\multicolumn{3}{|l|}{\textbf{Solo para transporte marítimo y por vías navegables}} \\ \hline
FAS & \textit{Free Alongside Ship} (franco al costado del buque) & El vendedor deja la
mercadería al costado del buque en el puerto de embarque. \\ \hline
FOB & \textit{Free On Board} (franco a bordo) & El vendedor entrega la mercadería a bordo
del buque en el puerto de embarque; desde ahí, costos y riesgos son del comprador. \\ \hline
CFR & \textit{Cost and Freight} (costo y flete) & El vendedor paga el flete marítimo hasta
el puerto de destino; el riesgo pasa al cargarla a bordo. \\ \hline
CIF & \textit{Cost, Insurance and Freight} (costo, seguro y flete) & Como CFR, pero el
vendedor contrata además un seguro de cobertura mínima. \\ \hline
\end{tabular}
\end{center}

\paragraph*{Errores comunes.}
\begin{itemize}
  \item \textbf{Usar FOB o CIF para envíos aéreos o por courier.} Son reglas marítimas. Para
        un envío aéreo de un instrumento o un repuesto corresponde FCA, CPT, CIP o DAP.
  \item \textbf{Usar siglas que ya no existen.} \textbf{DAF}, \textbf{DES}, \textbf{DEQ} y
        \textbf{DDU} se eliminaron en la versión 2010, y \textbf{DAT} fue reemplazado por
        DPU en 2020.
  \item \textbf{Confundir quién paga con quién corre el riesgo.} En CPT, CIP, CFR y CIF el
        vendedor paga el flete hasta destino, pero el riesgo pasa al comprador mucho antes,
        en origen.
\end{itemize}

\paragraph*{Ejemplo aeronáutico.}
Un taller de Quilmes compra a un proveedor de Miami un indicador de actitud reacondicionado.
\begin{itemize}
  \item Si compra \textbf{EXW Miami}, el taller tiene que retirar la pieza del depósito del
        proveedor y hacer también el trámite de exportación en Estados Unidos. Es la opción
        más barata en precio, pero la más complicada.
  \item Si compra \textbf{FCA Miami}, el proveedor entrega la pieza ya despachada de
        exportación al agente de carga o courier que eligió el taller. Es la regla más
        práctica para envíos aéreos.
  \item Si compra \textbf{CIP Ezeiza}, el proveedor paga el flete aéreo y el seguro hasta
        Ezeiza; el taller hace la importación y paga los tributos.
\end{itemize}

\subsubsection*{7.2. Importar y exportar en pequeña escala}
\begin{itemize}
  \item \textbf{Régimen de courier (pequeños envíos) para uso personal:} permite importar
        bienes para consumo personal con una franquicia de \textbf{USD 400 por envío, hasta
        5 envíos por año}; sobre esos primeros USD 400 se paga solo el IVA y, si
        corresponde, impuestos internos. El límite es de \textbf{USD 3.000 por envío}. Desde
        julio de 2026 las mismas condiciones se aplican a los envíos que entran por el
        correo oficial (puerta a puerta).
  \item \textbf{Exportación por courier:} desde 2026 ya no tiene límite de valor por envío,
        lo que facilita que pymes y emprendedores vendan al exterior por esta vía.
        A un cliente del exterior se le emite \textbf{factura E}.
  \item \textbf{Importación con fines comerciales:} para traer mercadería para revender o
        para producir hay que inscribirse como importador y operar por el régimen general,
        normalmente con un \textbf{despachante de aduana}. Recordar que quien importa bienes
        o servicios para revenderlos no puede estar en el monotributo.
\end{itemize}

% ------------------------------------------------------------
\section*{8. EL FINANCIAMIENTO: TIPOS DE PRÉSTAMOS BANCARIOS}

Un emprendimiento casi siempre necesita dinero que todavía no tiene: para comprar
herramientas, equipar un taller, comprar mercadería o pagar gastos mientras espera
cobrar. Un \textbf{préstamo} es un contrato por el cual una entidad entrega una suma de
dinero (el \textbf{capital}) y el tomador se compromete a devolverla en un plazo, más un
\textbf{interés}, que es el precio del dinero prestado.

\subsubsection*{8.1. Préstamos para personas}
\begin{itemize}
  \item \textbf{Personal:} de libre destino, sin garantía específica (``a sola firma''). Es
        rápido de obtener pero suele tener las tasas más altas. Los plazos habituales van
        de 6 a 72 meses.
  \item \textbf{Prendario:} para comprar un bien mueble registrable, generalmente un
        vehículo o una máquina. El bien queda \textbf{prendado} como garantía: si no se paga,
        el banco puede ejecutarlo.
  \item \textbf{Hipotecario:} para comprar, construir o ampliar una vivienda. El inmueble
        queda \textbf{hipotecado} como garantía. Son préstamos a largo plazo (hasta 20 o 30
        años).
  \item \textbf{UVA:} préstamos (personales o hipotecarios) cuyo capital se expresa en
        \textbf{Unidades de Valor Adquisitivo}, que se ajustan diariamente por inflación.
        La tasa de interés es baja, pero la cuota en pesos aumenta con la inflación.
  \item \textbf{Tarjeta de crédito y adelantos:} financiar el saldo de la tarjeta o pedir un
        adelanto en efectivo es una forma de préstamo, en general de costo muy alto.
\end{itemize}

\subsubsection*{8.2. Financiamiento para emprendimientos y empresas}
\begin{itemize}
  \item \textbf{Adelanto en cuenta corriente (descubierto):} el banco permite girar en
        descubierto hasta un límite acordado para cubrir faltantes de caja de corto plazo.
        Se pagan intereses solo por los días y el monto usados.
  \item \textbf{Descuento de cheques y e-cheqs:} el emprendedor que recibió un cheque de pago
        diferido (a 30, 60 o 90 días) lo entrega al banco y cobra hoy el importe, menos un
        descuento (interés y comisiones).
  \item \textbf{Factura de Crédito Electrónica MiPyME (FCEM):} cuando una pyme le vende a una
        gran empresa, emite este tipo de factura, que puede negociarse en bancos o en el
        mercado de capitales para cobrar antes del vencimiento.
  \item \textbf{Préstamo a sola firma o de capital de trabajo:} préstamos de corto plazo para
        comprar insumos o mercadería, pagar sueldos, etc.
  \item \textbf{Préstamo para inversión:} a mediano o largo plazo, para comprar maquinaria,
        equipos o instalaciones. Suele pedirse un proyecto o plan de negocios.
  \item \textbf{Leasing:} el banco compra el bien (por ejemplo, un banco de pruebas o un
        vehículo) y se lo alquila al emprendedor durante un plazo; al final, este puede
        comprarlo pagando un valor residual.
  \item \textbf{Garantías de SGR:} las Sociedades de Garantía Recíproca avalan a pymes que no
        tienen garantías propias, lo que facilita el acceso al crédito.
  \item \textbf{Microcréditos:} préstamos pequeños para microemprendedores, otorgados por
        bancos, organismos públicos u organizaciones sociales.
\end{itemize}

\subsubsection*{8.3. Cómo se mide el costo de un préstamo}
\begin{itemize}
  \item \textbf{TNA (Tasa Nominal Anual):} es la tasa ``de lista''. No refleja el costo real
        porque no considera la capitalización de los intereses.
  \item \textbf{TEA (Tasa Efectiva Anual):} considera la capitalización (por ejemplo, mensual).
        Siempre es mayor que la TNA.
  \item \textbf{CFT (Costo Financiero Total):} suma a la tasa todos los gastos: comisiones,
        seguros e impuestos. \textbf{Es el dato que hay que usar para comparar préstamos},
        siempre con igual monto, plazo y perfil del cliente.
  \item \textbf{Tasa fija o variable:} con tasa fija la cuota se conoce desde el principio;
        con tasa variable o ajuste UVA, la cuota puede cambiar.
\end{itemize}

\subsubsection*{8.4. Sistemas de amortización}
La cuota de un préstamo tiene dos partes: \textbf{interés} y \textbf{amortización}
(devolución del capital). Los dos sistemas más usados son:
\begin{itemize}
  \item \textbf{Sistema francés:} todas las cuotas son \textbf{iguales}. Al principio la cuota
        es casi todo interés y al final casi todo capital. Es el más común en préstamos
        personales.
  \item \textbf{Sistema alemán:} se devuelve la \textbf{misma porción de capital} en cada
        cuota y los intereses se calculan sobre el saldo, por lo que las cuotas son
        \textbf{decrecientes}. La primera cuota es más alta, pero se pagan menos intereses
        en total.
\end{itemize}

\paragraph*{Ejemplo.}
Préstamo de \$1.000.000 a 12 meses con una TNA del 60\% (5\% mensual):

\begin{center}
\small
\begin{tabular}{|l|r|r|}
\hline
 & \textbf{Sistema francés} & \textbf{Sistema alemán} \\ \hline
Cuota 1 & \$112.825,41 & \$133.333,33 \\ \hline
Cuota 12 & \$112.825,41 & \$87.500,00 \\ \hline
Total pagado & \$1.353.904,92 & \$1.325.000,00 \\ \hline
Total de intereses & \$353.904,92 & \$325.000,00 \\ \hline
\end{tabular}
\end{center}
La TEA de este préstamo es del 79,59\%, aunque la TNA anunciada sea del 60\%. Si además se
suman seguros y comisiones, el CFT será todavía mayor.

\subsubsection*{8.5. Derechos y recomendaciones}
\begin{itemize}
  \item Antes de firmar, revisar \textbf{cuánto dinero se recibe, cuánto se devuelve en
        total, el plazo, la cantidad de cuotas y el CFT}. La cuota no es el único dato que
        importa.
  \item No comprometer en cuotas más de lo que el emprendimiento puede pagar con sus
        ingresos reales.
  \item La \textbf{Central de Deudores del BCRA} registra los préstamos de los últimos 24
        meses y si se pagaron en término. Se puede consultar gratis en la web del BCRA
        (``situación crediticia''). Un atraso afecta el acceso a créditos futuros.
  \item Si hay un problema con un préstamo, primero se reclama a la entidad; si no responde
        o la respuesta no es satisfactoria, se puede reclamar ante el BCRA.
\end{itemize}

% ------------------------------------------------------------
\section*{9. EJERCICIO INTEGRADOR}

Un técnico abre un taller de reparación de equipos electrónicos y aviónica liviana, que
presta \textbf{servicios}. En los últimos 12 meses tuvo estos datos:
\begin{itemize}
  \item Facturación promedio: \$2.800.000 por mes (\$33.600.000 en el año).
  \item Superficie del taller: 50 m².
  \item Consumo eléctrico: 7.200 kWh en el año.
  \item Alquiler: \$250.000 por mes (\$3.000.000 en el año).
\end{itemize}

\paragraph*{Resolución.} Se busca la categoría que corresponde a cada parámetro en la tabla:
\begin{center}
\small
\begin{tabular}{|l|r|c|}
\hline
\textbf{Parámetro} & \textbf{Valor anual} & \textbf{Categoría} \\ \hline
Ingresos brutos & \$33.600.000 & E \\ \hline
Superficie & 50 m² & C \\ \hline
Energía eléctrica & 7.200 kWh & D \\ \hline
Alquileres & \$3.000.000 & C \\ \hline
\end{tabular}
\end{center}
Corresponde la categoría más alta de las cuatro: \textbf{categoría E}. Como presta servicios,
paga una cuota mensual de \textbf{\$119.811,45}. A sus clientes les emite \textbf{factura C}.

\paragraph*{Para resolver en clase.}
\begin{enumerate}
  \item ¿Qué categoría correspondería si, con los mismos datos, en lugar de reparar equipos
        vendiera repuestos? ¿Cuánto pagaría por mes?
  \item Si un cliente responsable inscripto le pide factura A, ¿puede emitirla? ¿Por qué?
  \item El taller necesita comprar un osciloscopio de \$3.000.000 y le ofrecen pagarlo en
        12 cuotas. Busquen en la web de dos bancos un préstamo para ese monto y comparen el
        CFT. ¿Cuál conviene?
  \item El taller compra por avión a un proveedor de Estados Unidos un repuesto de USD 1.200 y le
        cotizan ``FOB Miami''. ¿Qué error tiene esa cotización? ¿Qué Incoterm propondrían y por qué?
  \item Simulen en grupos la emisión de una factura C por un servicio de reparación,
        completando todos los datos que pide ``Comprobantes en línea''.
\end{enumerate}

\vspace{6pt}
{\footnotesize\textit{Fuentes: ARCA (arca.gob.ar: Monotributo – montos y categorías vigentes,
requisitos y exclusiones; Facturación – clases de comprobantes; Resoluciones Generales
5614/2024 y 5762/2025; Régimen de courier); ICC (Incoterms 2020); BCRA (Claves sobre préstamos, Central de Deudores). Datos consultados
el 25/11/2026.}}
